\begin{frame}[allowframebreaks]{Language}
    After the language is built
    
    \cite{rueda-2021-a-novel-wave-decomposition-for-oscillatory-signals} suggests using the superposition of two Frequency Modulated M\"obius (FMM) waves to represent two action potential waves.
    
    \begin{equation}
        \begin{aligned}
            V(t)
            =
            M
            &+
            A_1
            \cos
            \Biggl[
                \beta_1
                +
                2
                \arctan
                \Biggl(
                \omega_1
                \tan
                \Bigl(
                \frac{t-\alpha_1}{2}
                \Bigr)
                \Biggr)
                \Biggr]
            \\
            &+
            A_2
            \cos
            \Biggl[
                \beta_2
                +
                2
                \arctan
                \Biggl(
                \omega_2
                \tan
                \Bigl(
                \frac{t-\alpha_2}{2}
                \Bigr)
                \Biggr)
                \Biggr]
        \end{aligned}
    \end{equation}
    
    
    \section{Estimating one wave parameters}
        But the decoders have only learned to map one wave parameters so we should estimate the parameters of a single wave from the above somation using \textbf{derivative-free global optimization followed by local nonlinear least-squares refinement} to find the parameters on the left side of this equation
        \[
            \begin{aligned}
                &M
                +
                A\cos\left(
                         \beta
                         +
                         2\arctan\left(
                                     \omega\tan\left(\frac{t-\alpha}{2}\right)
                    \right)
                \right)
                \\
                &=
                \\
                &M_1
                +
                A_1\cos\left(
                           \beta_1
                    +
                    2\arctan\left(
                                \omega_1\tan\left(\frac{t-\alpha_1}{2}\right)
                    \right)
                \right)
                \\
                &+
                M_2
                +
                A_2\cos\left(
                           \beta_2
                    +
                    2\arctan\left(
                                \omega_2\tan\left(\frac{t-\alpha_2}{2}\right)
                    \right)
                \right)
            \end{aligned}
        \]
\end{frame}